%%
%% This is file `sample-sigconf.tex',
%% generated with the docstrip utility.
%%
%% The original source files were:
%%
%% samples.dtx  (with options: `sigconf')
%% 
%% IMPORTANT NOTICE:
%% 
%% For the copyright see the source file.
%% 
%% Any modified versions of this file must be renamed
%% with new filenames distinct from sample-sigconf.tex.
%% 
%% For distribution of the original source see the terms
%% for copying and modification in the file samples.dtx.
%% 
%% This generated file may be distributed as long as the
%% original source files, as listed above, are part of the
%% same distribution. (The sources need not necessarily be
%% in the same archive or directory.)
%%
%% The first command in your LaTeX source must be the \documentclass command.
\documentclass[sigplan,screen]{acmart}

\usepackage{code}
\usepackage{graphicx}
\usepackage{placeins}
\usepackage{tikz}
\usetikzlibrary{shapes, arrows, positioning}

%%
%% \BibTeX command to typeset BibTeX logo in the docs
\AtBeginDocument{%
  \providecommand\BibTeX{{%
    \normalfont B\kern-0.5em{\scshape i\kern-0.25em b}\kern-0.8em\TeX}}}

%% Rights management information.  This information is sent to you
%% when you complete the rights form.  These commands have SAMPLE
%% values in them; it is your responsibility as an author to replace
%% the commands and values with those provided to you when you
%% complete the rights form.


%%
%% Submission ID.
%% Use this when submitting an article to a sponsored event. You'll
%% receive a unique submission ID from the organizers
%% of the event, and this ID should be used as the parameter to this command.
%%\acmSubmissionID{123-A56-BU3}

%%
%% The majority of ACM publications use numbered citations and
%% references.  The command \citestyle{authoryear} switches to the
%% "author year" style.
%%
%% If you are preparing content for an event
%% sponsored by ACM SIGGRAPH, you must use the "author year" style of
%% citations and references.
%% Uncommenting
%% the next command will enable that style.
%%\citestyle{acmauthoryear}

%%% The following is specific to Onward! '24-ESSAYS and the paper
%%% '(Programs), Proofs and Refutations (and Tests and Mutants)'
%%% by Alex Groce.
%%%
\setcopyright{acmlicensed}
\acmDOI{10.1145/3689492.3689810}
\acmYear{2026}
\copyrightyear{2026}
\acmISBN{979-8-4007-1215-9/24/10}
\acmConference[Onward! '26]{Proceedings of the 2026 ACM SIGPLAN International Symposium on New Ideas, New Paradigms, and Reflections on Programming and Software}{October 3--9, 2026}{Oakland, CA, USA}
\acmBooktitle{Proceedings of the 2026 ACM SIGPLAN International
  Symposium on New Ideas, New Paradigms, and Reflections on
  Programming and Software (Onward! '26), October 3--9, 2026, Oakland, CA, USA}
\acmSubmissionID{onward26essays-p7-p}
%\received{2024-04-25}
%\received[accepted]{2024-08-08}

%%
%% end of the preamble, start of the body of the document source.
\begin{document}

%%
%% The "title" command has an optional parameter,
%% allowing the author to define a "short title" to be used in page headers.
\title{Another, Doubtless Very Different, Bill James: Baseball Numbers
and Software Numbers}

\author{Alex Groce}
\orcid{0000-0003-0273-4668}
\affiliation{%
  \institution{Northern Arizona University}
  \city{Flagstaff}
  \country{USA}
}
\email{agroce@gmail.com}

%% By default, the full list of authors will be used in the page
%% headers. Often, this list is too long, and will overlap
%% other information printed in the page headers. This command allows
%% the author to define a more concise list
%% of authors' names for this purpose.
\renewcommand{\shortauthors}{Alex Groce}

%%
%% The abstract is a short summary of the work to be presented in the
%% article.
\begin{CCSXML}
<ccs2012>
<concept>
<concept_id>10011007.10010940.10010992.10010998.10011001</concept_id>
<concept_desc>Software and its engineering~Dynamic analysis</concept_desc>
<concept_significance>500</concept_significance>
</concept>
<concept>
<concept_id>10011007.10011074.10011099.10011102.10011103</concept_id>
<concept_desc>Software and its engineering~Software testing and debugging</concept_desc>
<concept_significance>500</concept_significance>
</concept>
</ccs2012>
\end{CCSXML}

\ccsdesc[500]{Software and its engineering~Dynamic analysis}
\ccsdesc[500]{Software and its engineering~Software testing and debugging}

\keywords{software metrics}



\begin{abstract}
Software engineering metrics and baseball metrics have a lot in
common; this essay argues that we can make use of the common features,
with due attention to the critical differences, to guide future work on software
engineering metrics.
  \end{abstract}



\maketitle

\section{Introduction}

Harlan D. Mills (1919-1996) is the namesake of the Harlan D. Mills
Award given by the IEEE Computer Society to recognize “researchers and
practitioners who have demonstrated long-standing, sustained, and
meaningful contributions to the theory and practice of the information
sciences, focusing on contributions to the practice of software
engineering through the application of sound theory.”  First given in
1999, recipients include such luminaries as David Parnas (the first
winner), Elaine Weyuker, Bertrand Meyer, the Cousots, Gerard Holzmann,
Pamela Zave, Gail Murphy, Mark Harman, Matt Dwyer, David Harel, and,
most recently, Andreas Zeller.
Mills’ idea of the Chief Programmer Team influenced Brooks’ proposals
for how to avoid software disaster in \emph{The Mythical Man-Month}; his
Cleanroom conception informs every approach that is focused on
defect-prevention (from static analysis to aggressive unit testing to
mutation testing) rather than removing defects found in the field;
Mills was also a major force in introducing both statistics (viewing
tests as statistical experiments, which arguably also is a
foundational insight for all of random testing and fuzzing) and formal
specification (thus a contributor to the lines of thought that have
given the field Alloy, TLA+, Dafny, and other treasures).  In short,
Mills was a genuinely major figure in our field, though (other than having one
of the top awards in the field named after him) I think he’s one of
the least talked-about major visionaries we have. Indeed, one name for the Mills
Award, confusingly implied as being started in 1996, and perhaps by
ICSE/the ACM, in Wikipedia, is the “Harlan Mills Practical Visionary
Prize”.

I know the rough outlines of the above details because I am a software
engineering history buff.  But I also know that Mills was a command pilot for
the US Army Air Corps in WWII; that he rewrote the B-24 pilot training
program because he thought it was garbage; and that he worked at the
Institute for Advanced Study alongside Einstein, G\"odel, and von
Neumann. I know these things because Harlan Mills and his brother Eldon are the most
prominent characters in Chapter 4 of Alan Schwarz's book, \emph{The
  Numbers Game: Baseball's Lifelong Fascination with Statistics}~\cite{schwarz2004numbers}.
They appear there because, besides having a passion for applied mathematics,
operational research, and software engineering, Harlan Mills had a
passion for baseball.  In 1968, Harlan and Eldon Mills founded a
company called Computer Research in Sports.  Their goal was to go
beyond the statistics baseball
fans already knew to statistics that truly measured how a player
contributed to winning a game. For example, they wanted to look beyond
the much-beloved batting average, a statistic we will talk about
below that, one could argue, drove the country crazy a few times in the
years since baseball became America’s pastime.
Harlan and Eldon were among many high-powered
intellects (including multiple Nobel Prize winners and the popular
Harvard evolutionary-biology maverick, public intellectual, and superb
writer Stephen Jay Gould) pioneering an effort to truly understand
baseball, an effort that is documented in Schwarz’s engaging and in-depth look at
baseball’s love of statistics and how those statistics have developed
from the earliest, alien-looking, box-scores.

The connection between baseball statistics and software engineering,
however, goes beyond the curiosity that Mills was a pioneer in both
software engineering and baseball stats-savvy.  The obvious connection
is that the stats revolution in baseball was in substantial part tied
to the computer revolution; nothing like \url{https://www.retrosheet.org/}
could plausibly exist in a non-computerized world. This site,
which offers play-by-play accounts of over 100,000 pro baseball games
makes it possible to search for trends or individual facts, a task that
would be impossible even if a warehouse of filing cabinets had the
relevant data. Schwarz's book shows how computers were essential to
the growth of the modern baseball stats world, and takes us back to
when baseball work was done by “stealing” off-time on hulking company
and government mainframes. The deeper connection, however, is that baseball
was, for a very long time, tied to extremely traditional, and beloved,
statistics, but even stodgy, change-averse professional baseball front
and back offices eventually realized the grave limitations of these
statistics. Once they made this realization, they started making use of more sophisticated, nuanced, statistics
that truly reflected how to win the game.  That story:  famous,
widely-used numbers, lacking in proven operational relevance or at
least clearly sub-optimal in definition, should
sound familiar to students of software engineering.

\section{Build a Better Mousetrap/Metric}

%Baseball and software engineering have the following important similarities:
%data is easy to gather, hard to understand;
%a lot of people who like futzing with spreadsheets and
%coding up formulas are involved; and there’s both a lot of money and a lot of malarkey involved.
%Their differences are more stark:
%baseball is an extraordinarily structured thing, and the
%binary outcome of almost every event (did someone make an out? did
%someone score? did the team win?) is obvious; meanwhile in software
%engineering figuring out if a computer program works and was well-made
%turns out to be extremely annoyingly hard unless you wait twenty years
%and see if anyone is still using it and how much money it made.
%But
Perhaps we can quickly see what I mean by the move from traditional
stats to ones that really do \emph{mean more} by looking at a simple example from
software testing, paired with the most famous example from baseball
(and the only one I think I can explain tolerably well to software
people who don't know baseball).\footnote{We could also compare LOC produced to modern
software productivity measures that are less simpleminded, like DevOps' Change
Failure Rate~\cite{dora_cfr}.}

One way to determine if a program is well-tested is to take the test
suite for the program and look at the statement coverage: the percent
of all statements in the whole program that are executed at least once
by the test suite.  Obviously if a test suite fails to execute most of
the code in the program it is ostensibly testing, it is likely a
pretty bad test suite!  Similarly, if a test suite runs 90\% of the
code in the program, it’s presumably at least doing something
non-trivial right.

In baseball, for a very very long time, the king of statistics was
batting average.  Even if you know little
about baseball, even if you’re a European or from Mars \footnote{The best kid’s
book on baseball is Walter R. Brooks’ \emph{Freddy and the Baseball Team from Mars}, by the way).},
you probably know that in baseball, a batter is thrown balls by the
pitcher and tries to hit those balls, and then (when the ball is hit)
run to a series of three bases, before returning to the place from
which the ball is initially hit (home base).  Batting average looks at
how often a player hits the ball to advance to first, second, third,
or home base.  Batting average is not nonsense!  The most
sophisticated baseball statisticians tend to agree players celebrated
for their batting averages, such as Ty Cobb, Rogers Hornsby, and Ted
Williams, were superb, top-of-the-line players; the top players by
batting average now, like Aaron Judge, tend to be household names for
a good reason.  Batting average, like code
coverage, is a \emph{meaningful number}.

However, there are some problems with statement coverage, and code
coverage in general.  For one thing, maybe you only run a really
important line of code once in the whole test suite, and that in a
rather boring and uninteresting situation where the things the line of
code is supposed to do are not  seriously stressed (it’s a complex
multiplication, but one factor is zero, for example).  Or, worse yet,
you may run lots of code lots of times, but not check anything much
about whether the program did what it was supposed to do other than
“not crash.”  There are lots of ways a program can go wrong other than
outright crashing.  Mutation score~\cite{demillo1978hints,Discontents} asks “what do we really care about
in testing?  We want to find bugs.  Statement coverage is an awful
roundabout way to measure finding bugs, since bugs aren’t even
mentioned here.  Let’s see how a test suite does at finding bugs.”  In
mutation testing, a large number of essentially random bugs are
actually injected into a program, and the score is the percent of the
fake bugs the test suite was able to find.  It’s not perfect, but
it completely removes the two obvious problems
with statement coverage mentioned above, by focusing on \emph{finding
  the bug}: if you only run code in an uniteresting context, you won't
find a bug where the non-zero operand to multiplication is wrong; and
you won't find any bugs whose consequences you don't examine in your tests.

Similarly, the way batting average is defined, it doesn’t look at
every time a player comes up to home plate to try to hit the ball.
The numerator is “times the player hit the ball and got to a base” but
the denominator is “times the player came up to bat, minus times the
pitcher did such a bad job that the player was allowed to take a base
without hitting the ball, and the batter took proper advantage of this
by not swinging at bad pitches (or times the pitcher feared the player
so much this was intentionally done to keep the bat off the ball),
times the pitcher hit the player with the ball, times the catcher did
something weird, times the player intentionally hit the ball in a way
that tends not to let you reach base but has another
game-winning-relevant purpose, etc. etc. etc.).  On-base-percentage, a
revolution in the 1980s in understanding how to value hitters, counted
most everything, and looked (essentially) at “of times this player
came up to bat, how often did this player end up at least standing on
a base.”  Like mutation score, it focused on the \emph{more} essential (not
getting out before reaching a base, finding bugs) rather than the
easily and obviously measured (hitting the ball, running the code).
The typical non-baseball-fan probably still knows batting average, not
on-base-percentage; the typical software developer (who \emph{should}
be more like a baseball fan than someone who isn't sure if Babe Ruth
and Willie Mays played at the same time or not), alas, probably is
still more familiar with code coverage than mutation score, but the
game is changing, and highly visible adoption of mutation at, for example, Google~\cite{PetrovicMutationGoogle}
and Meta~\cite{BellerFacebookMutation} suggest a revolution in this
metric is underway.

\section{Why It Isn't Easy}

The above suggests that improving software engineering metrics is
easy:  just do what the baseball people did, and look at whether the
most widely used numbers are measuring the most important things.  

In fact, before we start on this, let's just start off with ways
software engineering is like baseball, and ways software engineering
is not like baseball.  Even if you don't know baseball, I am hoping
this list will be informative.

\subsection{Ways Software Engineering is Like Baseball}
\begin{itemize}
  \item Baseball and software engineering are both, when played for
    keeps by pros, cases where 1) it is fundamentally a team sport,
    and no one player, however good, can compete in the pros without a
    good team but 2)
    individual ability is extremely important, and the gap between the
    average player and the best is large and really matters.
    \item The ``game''  is deeply structured and appeals to
      stats-minded thinkers.
 
        \item Some of the best books in both fields are ``high end
          practitioners analyzing how they do their magic'':  for
          baseball, these range from Mathewson's legendary
          \emph{Pitching in a Pinch} to Hernandez's \emph {Pure
            Baseball}, with in betweens like Bob Gibson and Bob Gibson
          + Reggie Jackson (\emph{Pitch by Pitch} and \emph{Sixty Feet
            Six Inches}).  For software, we have legendary books like Brian Kernighan and
          Rob Pike's \emph{The Practice of Programming} and McConnell's
          \emph{Code Complete},  as well as \emph{The Pragmatic Programmer}, all
          the ``Uncle Bob'' \emph{Clean X} books, and so forth,
          including, arguably, much of Knuth's work.
          \item Crazy things that don't fit any model \emph{will}
            happen; you can understand and even in theory predict a
            phenomenon like the adoption of smarter metrics by
            management in baseball, or software companies adopting
            very good practices with respect to CI/testing; you can't
           predict a rare genius tech lead taking over a team and
           producing outsized results, or a couple of years of Billy
           Martin at the A's in the early 80s turning stolen bases and
           hit-and-run into remarkably good-for-the-talent seasons.
           Some things are not in any model.
         \end{itemize}


         
 None of these things contradict a narrative where ``doing what on-base-percentage
 does to batting average, but in more places'' is a fruitful and
 relatively simple path forwards for software metrics.  The problem is
 in the differences, and in fact in a few core differences between
 baseball and software, bringing us to:

 \subsection{Ways Software Engineering is Not Like Baseball}

 \begin{itemize}
   \item Baseball has a profoundly binary structure:  games are won or
     lost, and each game consists of a set of events with similarly
     clean, discrete outcomes:  an at-bat is generally a ``win''
     for the pitcher or
     the batter
     (perhaps with degrees of win for the batter);
     an inning has a score that is a set of of evaluations
    of scoring opportunities; even something as intrinsically complex
    and nuanced and hard-to-evaluate as fielding has obvious binary
    outcomes: was an error handed out?  was a potential double play
    converted? 
    \item Baseball's outcomes are almost entirely \emph{evaluated
        immediately}:  you do not need to wait six innings, much less
      a season, to know that a batter striking out is bad for the
      batter's team and good for the pitcher's team.  A design
      decision or API change often will not make its positive or
      negative payoffs clear for years or even decades.
      It's nonsensical to think of a baseball player who looks great
      on paper, racking up runs or outs for a team, but who in the
      long term has left such a load of baseball's version of
      technical debt that the player is a hidden liability, not an
      asset.
    \end{itemize}

 Notice that unlike with the similarities, I haven't bothered to
 discuss the ``software side'' of the issue, except implicitly:
 everyone likely to read this knows that most software actions and
 decisons don't have a simple, binary outcome, and that whatever
 outcome they do have is often greatly temporally delayed:  even if a
 line of code is simply and utterly wrong, the detection of that
 wrongness may wait hours, days, months, years, or even never come, if
 the bug is obscure and unlikely enough to be triggered.  If a guy in the outfield
 fails to catch a ball that should have easily been in his hands, it's
 clear; there are no subtle misses or catches.

 This presents a fundamental problem:  advanced statistics in baseball were, most
 famously, used to let a team without ``Yankees-money'' (for software-only
 people: read ``Google-money'') snap up undervalued players and
 perform tricks based on understanding of the deeper, better,
 statistics.  The \emph{Moneyball} movie, which stars Brad
 Pitt\footnote{That you will not see an equivalent software movie starring
   Brad Pitt is part of the point of this essay.}, has a famous bit of
 dialogue, where Billy Beane (the hero of the canonical ``applying
 advanced stats to actually winning ball games'' story), played by Pitt, talks to other
 people in the management structure of his team,
 who are trying to replace a certain critical individual player. Beane
 says this is a mistake.  He adds up the on-base-percentages of three players, and
 averages them, then says the goal needs to be not to replace any
 individual player, but to ``re-create him... in the aggregate'' and
 find three (ideally, under-valued) players whose on-base-percentage matches
 the three guys now missing from their team; not reproduce the
 expensive, hard-to-get, best of them, but match the average of the
 three, a much easier task.  In baseball, this pretty much works, if
 you know the values to use: using the highly-relevant
 on-base-percentage vs. much more famous, much  more lossy, batting
 average, matters here.  It's okay to take a guy who gets a lot of
 walks, for example: he's good at knowing when to hold off swinging at the ball,
 which is extremely useful for winning games; that this excites fans
 less than hitting the ball is ok, because winning games \emph{really}
 excites fans.

 In software, if you lose three developers, you cannot think about a
 formula for the ``good-code-percentage'' of the three, and try to
 match the mean value, with lesser-paid developers, none of whom match
 the star developer, but who, together, have the same
 rationally-measured contribution to software success.  Your departing
 developer has too few ``at-bats'' with too few binary ``good/bad''
 outcomes to turn into a real number; your possible replacements, not
 being highly (over-?) paid superstars, likely have even less reliable data;
 furthermore, knowing if any decisions any of these people have made
 were really good may take years; if you are a startup, you likely
 will be bought or have gone out of business before it's known if
 these players hit or struck out on their last at-bat, if it is ever
 known, which is unlikely.

The software game is just fundamentally hard to score; is Xerox's Palo
Alto team a winning or losing team?  They changed the world and gave
us modern software; they didn't make much money for Xerox.  There is
no baseball equivalent, really.  If you get a pennant or a World
Series title, you got it; nobody else gets to scoop up the rewards
of genius for what, on paper, is a losing season for you.

To emphasize this point, it's important to note that the most
important decisions in software are seldom individual lines-of-code
decisions, which map most naturally onto baseball's at-bats; the most
important decisions are design decisions, over-arching choices like
``do we use microservices or not?'' or ``what language should we write
this in?'' at the highest level, or ``should this class handle this
specialized scenario, or break it out into a separate class'' at a
lower level.  Baseball has important high-level decisions; a manager or
owner may prioritize hitting over pitching, power over speed, contact hitting
over power, etc., but these decisions still have quicker, clearer evaluation payoff, and
are more clearly constrained by budget or existing team, than
the most difficult, unconstrained, new-software-system architectural
decisions.  When you never make the playoffs, despite having a team of
very good players, in
five years, you can at least strongly suspect you chose wrong (and if
you are a manager, the owner will clarify things by firing you, probably); when you made an
architectural choice and you're winning over five years on deployment
flexibility, but debugging complexity is a nightmare costing you
money, reputation,
and burnt-out devs, you often can't
even guess if you made the right choice; it might well have gone worse
the other way.  

\section{How to Win the Game: Pick the Right Game to Play}

In baseball, if you have very bad pitchers, and no way to fix that,
you need to play the Boston Red Sox game of racking up a ton of runs
with an incredible offense; if you have great pitchers and no offense,
you need to play that game.  You may not have the ability to play the
ideal game (great pitching, great hitting, utter dominance) unless you
are the Yankees and have infinite money.

Similarly, if we software engineering thinkers want to up our metrics game, we
are going to have to look at where we can compete with baseball in
terms of developing real, advanced, meaningful statistics, and where
the structure of the game means we need to give up.

First, we have positive examples.  Mutation score is not a perfect way
to tell if a test suite is good or bad, but it is a much more powerful
approach to making a rational guess at that question than code
coverage alone, because it focuses on a genuinely meaningful binary
statistic in the same way as on-base-percentage:  does this test suite
detect this example of a bug?  This is repeated over a lot of at-bats,
and gives us a real measure of something important about a test suite
(the power it has to detect real, if hypothetical, bugs).

The above analysis suggests that trying to develop meaningful
statistics to evaluate individual players is, at minimum, going to be
vastly more difficult in software than in baseball.  Baseball has, in
fact, exploited the binary nature of a large number of immediately
evaluated events to produce sophisticated player statistics based on
concepts such as ``the delta between a given player and an aggregate
average, 'replacement' player at the same position'' (WAR, Wins Above
Replacement), runs created (a clever statistic invented by Bill James, about whom, more in
a little bit), an attempt to actually measure the genuine contribution
of an individual player to a team's overall set of winning games,
accounting for the player's individual role in play (win
shares, again a Bill James creation), and Pythagorean winning
percentage (again, James, and you mathy folks who don't know baseball can possibly even guess
what this one looks like).

The baseball analogy, though it does not give us a way to evaluate
individual developers well, does immediately expose some bad
metrics. Evaluating developers by lines of code is like evaluating
batters by how often they swing at a pitch. The pitch may have been a
ball; the swing, in itself, means nothing. What matters is getting on
base; in software, getting on base is writing useful, correct code that does not increase
technical debt or cognitive load (or at least does not increase it more than
is reasonable for the payoff). And often, one swing is better than
three; one line is better than five.

One might imagine that individual player evaluation must be left to
word-of-mouth, previous salary, resume scanning, and interviews
focusing on arbitrary and rather silly puzzles and coding questions,
but that perhaps software is doing a good job at project-level
estimation, where the hard-to-measure differences of individual
players are more evened-out.  Alas, even very well-intentioned
approaches such as COCOMO II~\cite{boehm2000cocomo} have a much higher error rate (30\% or
worse) than the current ability of Vegas oddsmakers or amateur
statisticians to predict the win-loss record of a team for an upcoming
season.  Large software projects remain notoriously over-budget and
over-schedule in a way dating back to Fred Brooks; baseball teams
sometimes surprise, but are fairly predictable in a way software
projects simply are not.

The way to play is to give up on the games you can't win right now.  In
the long run ``individual developers'' may in fact have reliable
statistics, if those developers are not humans with limited,
hard-to-evaluate, decisions but are LLMs that can be applied to
varieties of problems with known solutions for a huge number of
``at-bats'' and measured objectively on past problems, without undue
cost.  But for now, most developers are human, and there is no
shortcut to evaluating them; similarly, while some effort at project
budget and schedule estimation is necessary, it is unreasonable to
expect high precision much beyond what is obtained by simply asking expert
developers who are given a reasonable outline of the project scope.

But this does not mean that software has nowhere to go in improving
the quality of its statistics; it means that we have to look for
places where wins are possible, even if those are not the most
desirable statistical tools (developer evaluation tools for hiring and
firing,
project estimation tools for planning).  Testing, and more generally, CI, shows the way:
because testing has many binary outcomes immediately evaluated, 
successful/failed compilations, tests passed/failed, and static
analysis reports, it can be, and has been, turned into sophisticated
statistical measures that offer real power to smart managers and
developers.  Analysis in these settings has the ability to do real
historical examinations too, because old check-ins can be re-tested,
and you can actually ask if a proposed measure would have predicted
the eventual presence or absence of bugs in code~\cite{testedness}.

There are other places where such efforts seem like they might pay off:
\begin{itemize}
  \item Build and CI pipeline health in general (lots of discrete
  events, determination of flaky tests, etc.): real time series data
  \item Dependency and API stability:  when a library API breaks
    downstream consumers, that is often an observed event;
    vulnerabilities requiring critical patches are big-time, observed,
    ``strike-outs'' or ``errors''
    \item Code review turnaround (are bugs detected in code that
      passed code review, within a short time frame); this one can
      potentially point fingers at individual developers who may be good
      developers but are poor reviewers
    \item Incident/outage attribution (hard, the causality chains are a
      lot more opaque than the most complex baseball play, but not
      obviously hopeless)
      \item And, as noted, LLM/AI developers, who can be put through
        evaluation exercises that human developers cannot, without undue cost, undergo
      \end{itemize}

 These are simply the first few that come to my mind, and really I
 don't know anything about anything but testing, myself.

 \section{A New Bill James}

 In 2006, Time magazine named Bill James one of the 100
 most influential people in the world.  He worked for the Boston Red
 Sox for 17 years, during which they won four World Series
 championships.  He has undergraduate degrees in English and
 economics.  And, most importantly, by heroic effort, while a security
 guard at a pork-and-beans cannery, he asked a bunch of hard questions
 and gave data-based answers, focusing on things that could actually
 be answered given the data, and this work, after he wrote it up and
 fought his way through skeptical editors, kicked off a real
 revolution in baseball statistics, finally fulfilling the dream of
 Harlan Mills and other thinkers, by making baseball front offices turn their
 attention to numbers that meant something and were built for
 practical use.

 James' virtues, in baseball's ``numbers game,'' were a stubborn attention to the
 solvable problems at hand, a willingness to toil in obscurity until
 he had something worth saying, and
 the capability to write up his good ideas in a way that found a large
 audience, eventually including an audience consisting of the decision
 makers behind professional baseball teams.

 At the end of  his landmark book lamenting the collapse of moral
 philosophy as a vital program capable of directing human activity,
 the philosopher Alasdair MacIntyre wrote ``We are waiting not for a
 Godot, but for another---doubtless very different---St. Benedict.''
 Software engineering awaits another---doubtless very different---Bill James.

\bibliographystyle{ACM-Reference-Format}
\bibliography{bibliography}



\end{document}
